%%%%%%%%%%%%%%%%%%%%%%%%%%%%%%%%%%%%%%%%%%%%%%%%%%%%%%%%%%%%%%%%%%%%%%%%%%%%%%%
%                         CAPITULO 6
%           Protocolos de Comunicacion y Redes CAN
%%%%%%%%%%%%%%%%%%%%%%%%%%%%%%%%%%%%%%%%%%%%%%%%%%%%%%%%%%%%%%%%%%%%%%%%%%%%%%%

\chapter{Protocolos de Comunicacion y Redes CAN}

Hasta ahora hemos visto como se comunican dispositivos mediante UART, RS-232,
Wi-Fi, Bluetooth y redes TCP/IP. En el entorno industrial, sin embargo, existen
protocolos especificos diseñados para conectar sensores, actuadores, PLCs y
sistemas de control en tiempo real. Estos protocolos deben ser robustos, rapidos
y capaces de operar en entornos con ruido electrico y grandes distancias.

En este capitulo estudiamos los protocolos de comunicacion mas relevantes en
la industria: I2C, SPI, RS-485, Modbus y CAN. Cada uno tiene sus ventajas y
se emplea segun las necesidades de velocidad, distancia, numero de nodos y
fiabilidad del sistema.


%=============================================================================
\section{I2C: Inter-Integrated Circuit}
%=============================================================================

\subsection{Que es I2C}

El bus I2C (abreviatura de \textit{Inter-Integrated Circuits}) es un protocolo de
comunicacion serie sincrona desarrollado por Philips en 1982. Se diseño para
conectar circuitos integrados de baja velocidad dentro de un mismo dispositivo
o placa, como microcontroladores, sensores, memorias EEPROM y pantallas.

\begin{definicion}[I2C]
El \textbf{bus I2C} es un protocolo de comunicacion serie sincrono que utiliza
dos lineas para transmitir datos entre dispositivos: una linea de reloj (SCL) y
una linea de datos (SDA). Permite la conexion de multiples dispositivos en una
arquitectura de bus con multiples maestros y multiples esclavos.
\end{definicion}

La principal ventaja del I2C es su sencillez: solo necesita dos lineas de senal
mas masa (GND), lo que reduce el numero de pines necesarios en los
microcontroladores. Ademas, al ser un bus multi-maestro, cualquier dispositivo
conectado puede iniciar una transferencia si el bus esta libre.

\subsection{Senalizacion}

El bus I2C utiliza tres conexiones basicas:

\begin{itemize}
    \item \textbf{SCL} (\textit{System Clock}): linea de reloj que sincroniza la
    comunicacion entre todos los dispositivos. El maestro genera los pulsos de
    reloj.
    \item \textbf{SDA} (\textit{System Data}): linea de datos por la que se
    transfieren los bits de informacion entre maestro y esclavo.
    \item \textbf{GND} (Masa): referencia de tension comun para todos los
    dispositivos conectados al bus.
\end{itemize}

Ambas lineas (SCL y SDA) son \textit{open-drain} o \textit{open-collector}, lo que
significa que cada dispositivo solo puede poner la linea a nivel bajo (0) o
liberarla (1). Por eso se necesitan resistencias de pull-up externas
conectadas a la alimentacion (tipicamente 4.7 k$\Omega$ a 3.3 V o 5 V).

% TODO: imagen
\begin{figure}[H]
    \centering
    \fbox{\parbox{0.8\textwidth}{
        \centering
        \vspace{1cm}
        TODO: Esquema del bus I2C con maestro, dos esclavos y resistencias de pull-up.
        \vspace{1cm}
    }}
    \caption{Esquema basico del bus I2C con resistencias de pull-up.}
    \label{fig:i2c_esquema}
\end{figure}

\subsection{Transferencia de datos}

La comunicacion en I2C siempre la inicia un \textbf{maestro}. Cada esclavo tiene
una direccion unica de 7 bits (o 10 bits en modo extendido), que el maestro
envia al principio de cada transaccion.

\textbf{Escritura en un dispositivo esclavo:}

\begin{enumerate}
    \item El maestro envia una \textbf{condicion de inicio} (START): SDA pasa de
    alto a bajo mientras SCL permanece en alto.
    \item Envio de la \textbf{direccion del esclavo} (7 bits) mas el bit de
    lectura/escritura (0 = escritura).
    \item El esclavo responde con un bit de \textbf{reconocimiento} (ACK).
    \item El maestro envia la \textbf{direccion del registro interno} donde quiere
    escribir.
    \item El esclavo responde con ACK.
    \item El maestro envia el \textbf{byte de dato}.
    \item El esclavo responde con ACK.
    \item [Opcionalmente] se pueden enviar mas bytes de datos.
    \item El maestro envia una \textbf{condicion de parada} (STOP): SDA pasa de
    bajo a alto mientras SCL esta en alto.
\end{enumerate}

\textbf{Lectura desde un dispositivo esclavo:}

\begin{enumerate}
    \item Condicion de inicio (START).
    \item Direccion del esclavo con bit de escritura (0).
    \item ACK del esclavo.
    \item Envio de la direccion del registro interno a leer.
    \item ACK del esclavo.
    \item \textbf{Condicion de inicio reiterada} (Repeated START).
    \item Direccion del esclavo con bit de lectura (1).
    \item ACK del esclavo.
    \item El esclavo envia el \textbf{byte de dato}.
    \item El maestro responde con ACK (si quiere leer mas) o NACK (fin de lectura).
    \item Condicion de parada (STOP).
\end{enumerate}

\subsection{Ejemplo practico: sensor brujula digital CMPS03}

El sensor CMPS03 es una brujula digital que mide el campo magnetico terrestre
para determinar la orientacion. Una vez calibrado, ofrece una precision de 3 a 4
grados con resolucion de decimas.

Este sensor dispone de dos interfaces de salida:
\begin{itemize}
    \item \textbf{PWM} (modulacion en anchura de pulso): un pin genera pulsos cuya
    duracion es proporcional al angulo.
    \item \textbf{I2C}: permite leer el angulo directamente como un byte (0--255)
    o como dos bytes (0--3599, resolucion de 0.1 grados).
\end{itemize}

Para leer el angulo por I2C, el procedimiento es:
\begin{enumerate}
    \item Enviar START y la direccion del CMPS03 (0xC0) con bit de escritura.
    \item Enviar la direccion del registro interno 0x01 (angulo en 0--255).
    \item Enviar Repeated START.
    \item Enviar la direccion del CMPS03 (0xC1) con bit de lectura.
    \item Leer el byte de dato que contiene el angulo.
    \item Enviar STOP.
\end{enumerate}


%=============================================================================
\section{SPI: Serial Peripheral Interface}
%=============================================================================

\subsection{Que es SPI}

SPI es un protocolo de comunicacion serie sincrona desarrollado por Motorola en
1982. A diferencia de I2C, SPI esta pensado para alcanzar velocidades muy
superiores y opera en modo \textbf{full-duplex}, lo que significa que puede enviar
y recibir datos simultaneamente.

\begin{definicion}[SPI]
El \textbf{bus SPI} (\textit{Serial Peripheral Interface}) es un protocolo de
comunicacion serie sincrona de cuatro hilos que permite la transmision full-duplex
entre un dispositivo maestro y uno o varios esclavos. Utiliza lineas separadas
para reloj (SCK), datos de salida (MOSI), datos de entrada (MISO) y seleccion de
esclavo (SS).
\end{definicion}

SPI se utiliza habitualmente para comunicar un microcontrolador con perifericos
de alta velocidad como memorias Flash, pantallas TFT, convertidores ADC/DAC y
modulos de comunicacion.

\subsection{Senales del bus SPI}

Un bus SPI completo utiliza cuatro lineas:

\begin{itemize}
    \item \textbf{SCK} (\textit{Serial Clock}): senal de reloj generada por el
    maestro. Sincroniza la transmision de datos.
    \item \textbf{MOSI} (\textit{Master Out Slave In}): linea de datos del
    maestro al esclavo. El maestro envia informacion por esta linea.
    \item \textbf{MISO} (\textit{Master In Slave Out}): linea de datos del
    esclavo al maestro. El esclavo responde por esta linea.
    \item \textbf{SS} o \textbf{CS} (\textit{Slave Select} / \textit{Chip
    Select}): linea de control que el maestro utiliza para seleccionar a que
    esclavo se dirige. Normalmente activa en nivel bajo.
\end{itemize}

% TODO: imagen
\begin{figure}[H]
    \centering
    \fbox{\parbox{0.8\textwidth}{
        \centering
        \vspace{1cm}
        TODO: Esquema SPI con maestro, dos esclavos y lineas SCK, MOSI, MISO, SS.
        \vspace{1cm}
    }}
    \caption{Esquema del bus SPI con un maestro y dos esclavos.}
    \label{fig:spi_esquema}
\end{figure}

\subsection{Funcionamiento}

En SPI, el maestro siempre genera el reloj (SCK). En cada pulso de reloj, se
transfiere un bit en ambas direcciones simultaneamente: un bit sale por MOSI
y otro entra por MISO. Esto permite el modo full-duplex.

El maestro debe conocer de antemano cuantos bits espera recibir del esclavo,
porque debe mantener el reloj activo durante toda la transferencia. Si el esclavo
necesita tiempo para preparar la respuesta, el maestro puede introducir
pequenas pausas entre bytes.

\textbf{Seleccion de esclavo:}

Existen dos formas de conectar multiples esclavos:

\begin{enumerate}
    \item \textbf{SS independiente}: cada esclavo tiene su propia linea SS
    conectada al maestro. El maestro activa (baja) solo la linea del esclavo
    con el que quiere comunicarse.
    \item \textbf{Conexion en cascada (daisy chain)}: los esclavos se conectan
    en serie. El maestro envia datos a traves de todos los esclavos
    consecutivamente, como una cadena de registros de desplazamiento.
\end{enumerate}

\subsection{Comparacion: SPI vs I2C vs UART}

\begin{table}[H]
    \centering
    \caption{Comparativa de protocolos de comunicacion serie}
    \label{tab:comparativa_spi_i2c_uart}
    \small
    \begin{tabular}{|l|c|c|c|}
        \hline
        \textbf{Caracteristica} & \textbf{SPI} & \textbf{I2C} & \textbf{UART} \\
        \hline
        Tipo de senal & Sincrona & Sincrona & Asincrona \\
        \hline
        Modo duplex & Full-duplex & Half-duplex & Full-duplex \\
        \hline
        Numero de lineas & 4 (o mas) & 2 & 2 (TX/RX) \\
        \hline
        Velocidad tipica & Hasta 10+ Mbps & Hasta 3.4 Mbps & Hasta 1 Mbps \\
        \hline
        Maestros & Uno (tipico) & Multiples & Ninguno (peer-to-peer) \\
        \hline
        Esclavos & Multiples (por SS) & Multiples (por direccion) & 1 \\
        \hline
        Direccionamiento & Por linea SS & Por direccion 7/10 bits & Ninguno \\
        \hline
        Distancia & Corta (placa) & Corta (placa) & Media (15 m RS-232) \\
        \hline
    \end{tabular}
\end{table}

\textbf{Ventajas de SPI:}
\begin{itemize}
    \item Mayor velocidad que I2C y UART.
    \item Full-duplex: envia y recibe simultaneamente.
    \item El tamano de los mensajes puede ser arbitrario.
    \item Hardware sencillo y de bajo coste.
    \item Los esclavos no necesitan oscilador propio (el reloj lo genera el maestro).
\end{itemize}

\textbf{Desventajas de SPI:}
\begin{itemize}
    \item Necesita mas pines que I2C (especialmente si hay muchos esclavos).
    \item No hay confirmacion de recepcion por parte del esclavo (no hay ACK).
    \item Solo admite un maestro en la mayoria de implementaciones.
    \item Distancias muy cortas (dentro de la misma placa).
    \item El maestro debe conocer de antemano la longitud de la respuesta.
\end{itemize}


%=============================================================================
\section{RS-485: Comunicacion diferencial robusta}
%=============================================================================

\subsection{Introduccion}

RS-485 es un estandar de comunicacion serie muy utilizado en aplicaciones de
adquisicion de datos y control industrial. Su principal ventaja es la capacidad
de conectar multiples dispositivos en el mismo bus (hasta 32 transmisores y 32
receptores tipicamente) sobre distancias de hasta 1200 metros.

\begin{definicion}[RS-485]
El estandar \textbf{RS-485} (tambien conocido como EIA-485) define una interfaz
de comunicacion serie multipunto con transmision diferencial balanceada. Permite
la conexion de multiples transmisores y receptores en un bus semiduplex, siendo
altamente resistente al ruido electrico y apto para largas distancias.
\end{definicion}

\subsection{Caracteristicas principales}

Las caracteristicas mas destacadas de RS-485 son:

\begin{itemize}
    \item \textbf{Transmision diferencial}: los datos se envian mediante un par
    trenzado de hilos (A y B). Un hilo transmite la senal original y el otro su
    copia invertida. El receptor detecta la \textbf{diferencia de voltaje} entre
    ambos, lo que elimina el ruido comun que afecta por igual a ambos hilos.
    \item \textbf{Multipunto}: permite conectar hasta 32 dispositivos (transmisores
    y receptores) en el mismo bus. Con repetidores o transceivers avanzados se
    puede extender a mas nodos.
    \item \textbf{Semiduplex}: los dispositivos pueden enviar y recibir, pero no
    simultaneamente. En un instante dado, solo un equipo puede transmitir,
    mientras que todos los demas pueden recibir.
    \item \textbf{Inmunidad al ruido}: la transmision diferencial y el uso de par
    trenzado (a veces blindado) garantizan una alta resistencia a
    interferencias electromagneticas.
\end{itemize}

\subsection{Niveles de senal y cableado}

En RS-485, el estado logico se determina por la diferencia de voltaje entre las
lineas A y B:

\begin{itemize}
    \item \textbf{Bit 1 (Marca / recessive)}: $V_A - V_B < -200$ mV.
    \item \textbf{Bit 0 (Espacio / dominant)}: $V_A - V_B > +200$ mV.
\end{itemize}

El cableado basico consiste en un par trenzado para los datos (A y B), mas
lineas opcionales de alimentacion (0 V y 5 V) para alimentar dispositivos del
bus. Es recomendable colocar \textbf{terminaciones de linea} (resistencias de 120
$\Omega$) al principio y al final del bus para evitar reflexiones de senal.

% TODO: imagen
\begin{figure}[H]
    \centering
    \fbox{\parbox{0.8\textwidth}{
        \centering
        \vspace{1cm}
        TODO: Esquema de bus RS-485 con par trenzado, terminaciones de 120 Ohm
        y varios nodos conectados.
        \vspace{1cm}
    }}
    \caption{Topologia tipica de un bus RS-485 con terminaciones.}
    \label{fig:rs485_topologia}
\end{figure}

\subsection{Distancia vs velocidad}

Una regla practica en RS-485 es que el producto de la longitud del cable (en
metros) por la velocidad de datos (en bits por segundo) no debe superar
$10^8$:

\begin{equation}
    \mathrm{Longitud\,(m)} \times \mathrm{Velocidad\,(bps)} \leq 10^8
\end{equation}

\begin{table}[H]
    \centering
    \caption{Relacion entre distancia y velocidad en RS-485}
    \label{tab:rs485_distancia}
    \begin{tabular}{|c|c|}
        \hline
        \textbf{Distancia} & \textbf{Velocidad maxima aproximada} \\
        \hline
        15 m & 10 Mbps \\
        \hline
        100 m & 1 Mbps \\
        \hline
        500 m & 200 kbps \\
        \hline
        1200 m & 100 kbps \\
        \hline
    \end{tabular}
\end{table}

\subsection{RS-232 vs RS-485}

\begin{table}[H]
    \centering
    \caption{Comparativa entre RS-232 y RS-485}
    \label{tab:rs232_vs_rs485}
    \begin{tabular}{|l|c|c|}
        \hline
        \textbf{Caracteristica} & \textbf{RS-232} & \textbf{RS-485} \\
        \hline
        Tipo de comunicacion & Punto a punto & Multipunto \\
        \hline
        Duplex & Full-duplex & Semiduplex \\
        \hline
        Tipo de senal & Desbalanceada (unipolar) & Balanceada (diferencial) \\
        \hline
        Niveles de voltaje & $\pm 3$ V a $\pm 15$ V & Diferencial 200 mV \\
        \hline
        Max. dispositivos & 1 TX + 1 RX & 32 TX + 32 RX (estandar) \\
        \hline
        Velocidad maxima & 1 Mbps a 15 m & 10 Mbps a 15 m \\
        \hline
        Distancia maxima & $\sim$ 15 m & $\sim$ 1200 m a 100 kbps \\
        \hline
        Inmunidad al ruido & Baja & Alta \\
        \hline
        Lineas de control & DTR, DSR, RTS, CTS & Ninguna (solo datos) \\
        \hline
    \end{tabular}
\end{table}


%=============================================================================
\section{Modbus: protocolo de comunicacion industrial}
%=============================================================================

\subsection{Que es Modbus}

Modbus es uno de los protocolos de comunicacion mas extendidos en la
automatizacion industrial. Fue desarrollado en 1979 por Modicon (hoy Schneider
Electric) para su gama de PLCs. Hoy en dia es un estandar abierto que permite
la comunicacion entre controladores, sensores, actuadores y sistemas SCADA.

\begin{definicion}[Modbus]
El protocolo \textbf{Modbus} es un protocolo de comunicacion abierto para la
transmision de datos entre dispositivos electronicos en entornos industriales.
Se situa en los niveles 1, 2 y 7 del modelo OSI. Funciona sobre arquitectura
maestro/esclavo (en RTU) o cliente/servidor (en TCP/IP).
\end{definicion}

Modbus es independiente del medio fisico: puede funcionar sobre RS-232,
RS-485, Ethernet TCP/IP, o incluso por radio. Esto lo hace extremadamente
versatil.

\subsection{Variantes de Modbus}

A lo largo de los años han aparecido diferentes variantes de Modbus adaptadas a
diversos medios fisicos:

\begin{table}[H]
    \centering
    \caption{Variantes del protocolo Modbus}
    \label{tab:modbus_variantes}
    \begin{tabular}{|l|p{8cm}|}
        \hline
        \textbf{Variante} & \textbf{Descripcion} \\
        \hline
        Modbus RTU & Implementacion mas comun. Usa comunicacion serie (RS-232/RS-485)
        con representacion binaria compacta de los datos. Requiere CRC para
        verificacion. \\
        \hline
        Modbus ASCII & Usa comunicacion serie con caracteres ASCII. Mas lento que
        RTU. Utiliza checksum LRC. \\
        \hline
        Modbus TCP/IP & Funciona sobre Ethernet TCP/IP, puerto 502. No requiere
        checksum adicional porque TCP ya lo proporciona. \\
        \hline
        Modbus RTU sobre IP & Similar a Modbus TCP pero incluye CRC en la carga
        util, como en RTU. \\
        \hline
        Modbus UDP & Variante experimental sobre UDP para reducir overhead. \\
        \hline
        Modbus Plus (MB+) & Version propietaria de Schneider Electric. Soporta
        comunicacion peer-to-peer entre multiples maestros. Requiere hardware
        especial (co-procesador HDLC). \\
        \hline
    \end{tabular}
\end{table}

\subsection{Modbus RTU en detalle}

Modbus RTU (\textit{Remote Terminal Unit}) es la variante mas utilizada en el
entorno industrial. Se basa en una arquitectura \textbf{maestro-esclavo} sobre
un bus serie (generalmente RS-485).

\textbf{Estructura de la trama RTU:}

\begin{table}[H]
    \centering
    \caption{Formato de trama Modbus RTU}
    \label{tab:modbus_trama}
    \begin{tabular}{|l|c|p{6cm}|}
        \hline
        \textbf{Campo} & \textbf{Tamano} & \textbf{Descripcion} \\
        \hline
        Direccion del esclavo & 1 byte & Identifica el dispositivo destino (1--247) \\
        \hline
        Codigo de funcion & 1 byte & Indica la operacion a realizar (leer, escribir, etc.) \\
        \hline
        Datos & N bytes & Parametros de la funcion y valores \\
        \hline
        CRC & 2 bytes & Codigo de redundancia ciclica para deteccion de errores \\
        \hline
    \end{tabular}
\end{table}

El dispositivo \textbf{maestro} (por ejemplo un PLC o sistema SCADA) inicia la
comunicacion enviando una trama al esclavo. Solo el esclavo con la direccion
indicada responde. Los demas escuchan pero ignoran el mensaje.

\textbf{Codigos de funcion mas comunes:}
\begin{itemize}
    \item \textbf{01}: Leer bobinas (coils) -- salidas digitales.
    \item \textbf{02}: Leer entradas digitales.
    \item \textbf{03}: Leer registros de retencion (holding registers) -- 16 bits.
    \item \textbf{04}: Leer registros de entrada (input registers).
    \item \textbf{05}: Escribir una bobina.
    \item \textbf{06}: Escribir un registro de retencion.
    \item \textbf{15}: Escribir multiples bobinas.
    \item \textbf{16}: Escribir multiples registros.
\end{itemize}

\subsection{RS-232 vs RS-485 en Modbus}

La eleccion entre RS-232 y RS-485 depende de las necesidades de la aplicacion:

\begin{itemize}
    \item \textbf{RS-232}: adecuado para conexiones punto a punto a corta distancia
    (hasta 15 m). Usa conectores DB9 o DB25. Ideal para conectar un PLC
    directamente a un panel HMI o a un PC de programacion.
    \item \textbf{RS-485}: emplea un esquema de bus multidrop que permite conectar
    multiples dispositivos en la misma linea. Ideal para distancias largas
    (hasta 1200 m) y entornos con ruido electrico. Es el estandar de facto en
    redes de campo industriales.
\end{itemize}


%=============================================================================
\section{CAN: Controller Area Network}
%=============================================================================

\subsection{Introduccion}

CAN (\textit{Controller Area Network}) es un bus de comunicaciones serie
diseñado originalmente por Bosch en 1986 para la industria automotriz. Hoy en
dia es un estandar internacional (ISO 11898) y se emplea ampliamente en
automocion, maquinaria industrial, robots, sistemas medicos y avionica.

\begin{definicion}[CAN]
El \textbf{bus CAN} (\textit{Controller Area Network}) es un protocolo de
comunicacion de datos en serie de alta integridad para aplicaciones en tiempo
real. Utiliza un bus de dos hilos con arbitraje basado en prioridad, difusion de
mensajes, y excelentes capacidades de deteccion de errores y confinamiento de
fallos.
\end{definicion}

CAN se distingue de otros protocolos porque no utiliza direcciones de nodo.
En su lugar, cada mensaje lleva un \textbf{identificador} que define tanto la
prioridad como el contenido del mensaje. Cualquier nodo puede enviar un
mensaje cuando el bus esta libre, y todos los demas nodos reciben ese mensaje,
decidiendo individualmente si les interesa o no.

\subsection{Caracteristicas principales}

Las caracteristicas mas destacadas de CAN son:

\begin{itemize}
    \item \textbf{Acceso multi-maestro}: cualquier nodo puede iniciar la
    transmision cuando el bus esta libre.
    \item \textbf{Difusion de mensajes} (broadcast): todos los nodos reciben cada
    mensaje. No hay comunicacion punto a punto.
    \item \textbf{Prioridad de mensajes}: el identificador determina la prioridad.
    Los mensajes con identificador mas bajo tienen mayor prioridad.
    \item \textbf{Longitud limitada}: hasta 8 bytes de datos por mensaje.
    \item \textbf{Velocidad}: hasta 1 Mbit/s en distancias cortas (hasta 40 m a 1
    Mbps; 125 kbps a 500 m aproximadamente).
    \item \textbf{Deteccion de errores}: mecanismos de supervision de bits,
    CRC, reconocimiento y relleno de bits garantizan una tasa de error
    residual extremadamente baja.
\end{itemize}

\subsection{Tipos de tramas CAN}

En CAN existen cuatro tipos de tramas:

\begin{enumerate}
    \item \textbf{Trama de datos}: transmite un mensaje al bus. Contiene el
    identificador, los datos (0--8 bytes) y campos de control y CRC.
    \item \textbf{Trama remota}: solicita la transmision de un mensaje con un
    identificador especifico. Un nodo que produzca ese mensaje respondera
    enviando la trama de datos correspondiente.
    \item \textbf{Trama de error}: cualquier nodo que detecte una condicion de
    error envia una trama de error para notificar a todos los demas nodos.
    \item \textbf{Trama de sobrecarga}: similar a la trama de error, pero indica
    que un nodo no esta listo para recibir mas datos y necesita un retraso.
\end{enumerate}

\subsection{Formato de trama de datos}

Una trama de datos CAN (formato estandar con identificador de 11 bits) contiene
los siguientes campos:

\begin{table}[H]
    \centering
    \caption{Campos de una trama de datos CAN (formato estandar)}
    \label{tab:can_trama}
    \begin{tabular}{|l|c|p{6cm}|}
        \hline
        \textbf{Campo} & \textbf{Bits} & \textbf{Descripcion} \\
        \hline
        Inicio de trama (SOF) & 1 & Bit dominante que marca el inicio \\
        \hline
        Identificador + RTR & 12 & 11 bits de ID + 1 bit RTR (Remote Transmission Request) \\
        \hline
        Control & 6 & Indica tamano de datos y formato (estandar/extendido) \\
        \hline
        Datos & 0--64 & De 0 a 8 bytes de informacion util \\
        \hline
        CRC & 15 + 1 & Codigo de redundancia ciclica + delimitador \\
        \hline
        ACK & 2 & Ventana de reconocimiento por parte de los receptores \\
        \hline
        Fin de trama (EOF) & 7 & Bits recesivos que marcan el final \\
        \hline
        Intermission (IFS) & 3 & Tiempo minimo entre tramas \\
        \hline
    \end{tabular}
\end{table}

\textbf{Formato estandar vs extendido:}

\begin{itemize}
    \item \textbf{CAN estandar} (CAN 2.0A): identificador de 11 bits (2048
    mensajes posibles, aunque algunos IDs estan reservados).
    \item \textbf{CAN extendido} (CAN 2.0B): identificador de 29 bits (mas de 500
    millones de mensajes posibles).
\end{itemize}

Ambos formatos pueden coexistir en el mismo bus. La distincion se realiza
mediante el bit \textbf{IDE} (Identifier Extension).

\subsection{Arbitraje de bus: dominante y recesivo}

CAN utiliza un mecanismo de arbitraje \textbf{no destructivo} basado en los
niveles logicos del bus:

\begin{itemize}
    \item \textbf{Nivel dominante} (0 logico): el estado que impone su valor
    cuando un nodo lo transmite.
    \item \textbf{Nivel recesivo} (1 logico): el estado que solo se mantiene si
    \textit{ningun} nodo transmite un dominante.
\end{itemize}

El bus CAN es como una \textbf{AND cableada}: si un nodo pone dominante y otro
recesivo, el resultado es dominante. Esto permite que varios nodos comiencen a
transmitir simultaneamente, pero solo el que envie el mensaje con el
\textbf{identificador mas bajo} (mas dominantes en los primeros bits) gane el
arbitraje.

\textbf{Proceso de arbitraje:}
\begin{enumerate}
    \item Dos o mas nodos detectan que el bus esta libre y comienzan a transmitir
    al mismo tiempo.
    \item Cada nodo envia su identificador bit a bit.
    \item Cada nodo transmisor supervisa el nivel real del bus.
    \item Si un nodo envia un recesivo (1) pero ve un dominante (0) en el bus,
    significa que otro nodo tiene un identificador con mayor prioridad (ID mas
    bajo). Este nodo se retira y pasa a modo receptor.
    \item El nodo con el ID mas bajo gana el arbitraje y continua transmitiendo
    sin perdida de tiempo ni datos.
    \item Los nodos que perdieron reintentaran automaticamente cuando el bus vuelva
    a estar libre.
\end{enumerate}

% TODO: imagen
\begin{figure}[H]
    \centering
    \fbox{\parbox{0.8\textwidth}{
        \centering
        \vspace{1cm}
        TODO: Diagrama temporal de arbitraje CAN con dos nodos compitiendo.
        Nodo A (ID 0x123) gana sobre Nodo B (ID 0x456).
        \vspace{1cm}
    }}
    \caption{Arbitraje no destructivo en el bus CAN.}
    \label{fig:can_arbitraje}
\end{figure}

\subsection{Supervision de bits y relleno de bits}

\textbf{Supervision de bits (bit monitoring):}
Cada nodo transmisor supervisa el nivel del bus bit a bit y lo compara con lo
que el mismo esta enviando. Si detecta una discrepancia, sabe que ha ocurrido
un error (o que perdio el arbitraje).

\textbf{Relleno de bits (bit stuffing):}
Para garantizar suficientes transiciones en el bus (necesarias para la
sincronizacion), CAN inserta un bit de polaridad opuesta despues de 5 bits
consecutivos del mismo valor. Este bit de relleno es transparente para la capa
superior y se elimina en recepcion.

La longitud de la trama varia debido al relleno. En el peor caso (identificador
de 11 bits, datos maximos):

\begin{table}[H]
    \centering
    \caption{Longitud de trama CAN con relleno de bits}
    \label{tab:can_longitud}
    \begin{tabular}{|c|c|c|c|}
        \hline
        \textbf{Datos} & \textbf{Sin relleno} & \textbf{Promedio} & \textbf{Peor caso} \\
        \hline
        0 bytes & 47 bits & 49 bits & 55 bits \\
        \hline
        8 bytes & 111 bits & 114 bits & 135 bits \\
        \hline
    \end{tabular}
\end{table}

\subsection{Errores y confinamiento de fallos}

CAN dispone de sofisticados mecanismos para detectar y gestionar errores:

\textbf{Tipos de errores:}
\begin{enumerate}
    \item \textbf{Error de bit}: un transmisor envia un bit pero lee el nivel
    opuesto en el bus.
    \item \textbf{Error de relleno}: se detectan 6 bits consecutivos del mismo
    valor.
    \item \textbf{Error CRC}: el calculo CRC del receptor no coincide con el
    campo CRC recibido.
    \item \textbf{Error de formato}: bits de formato fijo en posiciones incorrectas.
    \item \textbf{Error de reconocimiento (ACK)}: ningun receptor envio un ACK
    dominante.
\end{enumerate}

\textbf{Confinamiento de fallos (fault confinement):}
Cada nodo CAN mantiene dos contadores de errores:
\begin{itemize}
    \item \textbf{TEC} (\textit{Transmit Error Counter}): cuenta errores en
    transmision.
    \item \textbf{REC} (\textit{Receive Error Counter}): cuenta errores en
    recepcion.
\end{itemize}

Segun el valor de estos contadores, un nodo puede estar en tres estados:
\begin{enumerate}
    \item \textbf{Activo} (Error Active): estado normal. El nodo puede enviar
    tramas de error activas (6 bits dominantes).
    \item \textbf{Pasivo} (Error Passive): cuando un contador supera 127. El nodo
    solo puede enviar tramas de error pasivas (6 bits recesivos) y debe esperar
    8 bits adicionales antes de retransmitir.
    \item \textbf{Bus Off}: cuando TEC supera 255. El nodo se desconecta del bus
    y no puede transmitir ni recibir hasta que el problema se resuelva.
\end{enumerate}

Este sistema de \textbf{voto mayoritario} asegura que un nodo defectuoso no
sature el bus con tramas de error, protegiendo al resto de la red.

\subsection{Topologia y cableado}

CAN utiliza una \textbf{topologia de bus} con dos lineas principales:

\begin{itemize}
    \item \textbf{CAN\_H}: linea de bus en estado dominante alto.
    \item \textbf{CAN\_L}: linea de bus en estado dominante bajo.
\end{itemize}

El bus debe terminarse en ambos extremos con resistencias de 120 $\Omega$ para
evitar reflexiones. En topologias cortas, una sola terminacion puede bastar.

% TODO: imagen
\begin{figure}[H]
    \centering
    \fbox{\parbox{0.8\textwidth}{
        \centering
        \vspace{1cm}
        TODO: Topologia de bus CAN con terminaciones, nodos conectados y
        conector DB9 con pines CAN\_H, CAN\_L, CAN\_GND y CAN\_SHLD.
        \vspace{1cm}
    }}
    \caption{Topologia de bus CAN segun la norma ISO 11898.}
    \label{fig:can_topologia}
\end{figure}

\textbf{Conector DB9 comun en CAN:}

\begin{table}[H]
    \centering
    \caption{Asignacion de pines en conector DB9 para CAN}
    \label{tab:can_db9}
    \begin{tabular}{|c|l|l|}
        \hline
        \textbf{Pin} & \textbf{Senal} & \textbf{Descripcion} \\
        \hline
        2 & CAN\_L & Linea de bus CAN (bajo dominante) \\
        \hline
        3 & CAN\_GND & Masa de referencia CAN \\
        \hline
        5 & CAN\_SHLD & Blindaje opcional \\
        \hline
        7 & CAN\_H & Linea de bus CAN (alto dominante) \\
        \hline
        9 & CAN\_V+ & Alimentacion externa opcional \\
        \hline
    \end{tabular}
\end{table}

\subsection{Controladores CAN: Basic vs Full CAN}

Los chips controladores CAN se clasifican en dos categorias:

\begin{itemize}
    \item \textbf{Basic CAN}: ofrece un buffer de recepcion FIFO y un buffer de
    transmision. Es una solucion de bajo coste, pero requiere que el CPU atienda
    rapidamente los mensajes para no perder datos en redes con mucho trafico.
    \item \textbf{Full CAN}: proporciona multiples buffers de mensaje programables
    (tanto para recepcion como para transmision). El controlador puede filtrar
    mensajes por identificador y almacenarlos en buffers dedicados sin
    intervencion del CPU. Permite gestionar alto trafico de CAN con bajo
    rendimiento de CPU.
\end{itemize}

La mayoria de los microcontroladores modernos integran un controlador CAN Full
o al menos Basic CAN, lo que facilita su uso en aplicaciones industriales sin
necesidad de hardware externo.

\subsection{Resumen: SPI, I2C, RS-485, Modbus y CAN}

\begin{table}[H]
    \centering
    \caption{Resumen comparativo de protocolos industriales}
    \label{tab:resumen_protocolos}
    \small
    \begin{tabular}{|l|c|c|c|c|c|}
        \hline
        \textbf{Caracteristica} & \textbf{I2C} & \textbf{SPI} & \textbf{RS-485} & \textbf{Modbus} & \textbf{CAN} \\
        \hline
        Tipo & Sincrono & Sincrono & Diferencial & Protocolo & Bus serie \\
        \hline
        Duplex & Half & Full & Half & Half/Full & Half \\
        \hline
        Nodos & 127 (7 bits) & Limitado por SS & 32+ & 247 (RTU) & Ilimitado \\
        \hline
        Distancia & Corta (placa) & Corta (placa) & 1200 m & 1200 m & 500 m+ \\
        \hline
        Velocidad & 3.4 Mbps & 10+ Mbps & 10 Mbps & 115.2 kbps & 1 Mbps \\
        \hline
        Arbitraje & Software & SS hardware & Maestro unico & Master/slave & No destructivo \\
        \hline
        Deteccion errores & ACK & Ninguna & CRC/paridad & CRC (RTU) & CRC + bit stuff \\
        \hline
        Uso tipico & Sensores, EEPROM & Pantallas, Flash & Redes de campo & PLCs, SCADA & Automocion, robots \\
        \hline
    \end{tabular}
\end{table}


%=============================================================================
\section{Resumen de la Unidad 6}
%=============================================================================

En este capitulo hemos estudiado los protocolos de comunicacion mas
relevantes para la interconexion de dispositivos en entornos industriales:

\begin{enumerate}
    \item \textbf{I2C} es un bus sincrono de dos hilos (SDA, SCL) que permite
    conectar multiples maestros y esclavos en cortas distancias. Es ideal para
    sensores y circuitos integrados dentro de una placa.

    \item \textbf{SPI} es un bus sincrono de cuatro hilos (SCK, MOSI, MISO, SS)
    que opera en full-duplex a muy alta velocidad. Perfecto para perifericos
    rapidos como memorias Flash o pantallas.

    \item \textbf{RS-485} es un estandar electrico de transmision diferencial
    balanceada que permite conectar hasta 32 dispositivos en un bus de hasta
    1200 m. Es la base fisica de muchas redes industriales.

    \item \textbf{Modbus} es un protocolo de aplicacion abierto muy utilizado en
    automatizacion. Funciona sobre RS-232, RS-485 o TCP/IP, con arquitectura
    maestro/esclavo en RTU y cliente/servidor en TCP.

    \item \textbf{CAN} es un bus de comunicacion de alta fiabilidad con
    arbitraje no destructivo por prioridad, difusion de mensajes, excelente
    deteccion de errores y confinamiento de fallos. Se usa en automocion,
    maquinaria industrial, robots y sistemas medicos.
\end{enumerate}

La eleccion del protocolo adecuado depende del entorno: para sensores dentro de
una placa, I2C o SPI; para redes de campo en una fabrica, RS-485 con Modbus o
CAN; para sistemas criticos en tiempo real, CAN es la opcion mas robusta.
